% !TEX root = ../../cookbook.tex
\newpage
\recipe{Pasta alle acciughe}{Ricetta originale: \url{ricette.giallozafferano.it}. Per 4 persone}{}{}

\subsection*{Ingredients}
\begin{ingredients}
	\item 400 g di pasta (idealmente spaghetti)
	\item 30 g di acciughe sott'olio sgocciolate
	\item 35 g di concentrato di pomodoro
	\item 30 g di pinoli
	\item 1 spicchio d'aglio
	\item 1 ciuffetto di prezzemolo
	\item \sfrac{1}{2} peperoncino fresco
	\item Vino bianco per sfumare
	\item Olio extravergine d'oliva q.b.
	\item 1 cucchiaio da minestra di mascarpone per persona (facoltativo)
\end{ingredients}

% --- Preparation ---
Mettere a bollire la \textbf{pasta}, tenendo da parte un bicchierino dell'acqua di cottura.
Dividere le \textbf{acciughe} in terzi, versandone l'olio direttamente in una padella da scaldare. 
Sciogliere \sfrac{1}{3} delle acciughe nella padella a fuoco medio. 
Intanto, su un'asse, tritare insieme il resto delle acciughe con i \textbf{pinoli}, il \textbf{prezzemolo}, l'\textbf{aglio}, e il \textbf{peperoncino}.

Quando le acciughe in padella saranno sciolte, aggiungere un po' di acqua di cottura e il \textbf{concentrato di pomodoro}.
Lasciar ridurre, quindi aggiungere il trito aromatico e soffriggerlo per un paio di minuti\footnote{Se necessario, aggiungere un po' più d'olio d'oliva.}. 
Sfumare col \textbf{vino bianco}. 
Aggiungere la pasta scolata al dente alla padella, eventualmente con un altro po' dell'acqua di cottura. 
Se piace, aggiungere un cucchiaio di mascarpone fresco per persona. 

% --- Description ---
\begin{center}
	\begin{minipage}{\textwidth} % Expanded to full text width for 3 columns
		\begin{recipeblock}
			\begin{multicols}{3}
				\footnotesize
				\centering % Centers text inside each column
				\setlength{\parskip}{1em} % Adds space between stanzas
				
				Acciughe \\
				Acciughe salate. \\
				Acciughe sebbele \\
				Trote salmonate. 
				
				Salmone \\
				Capparossoli,\\
				Risotto \\
				Al nero di seppia. 
				
				Spaghetti \\
				All'aragosta. 
				
				Anguilla \\
				Frittura mista, \\
				Polpa di granchio \\
				Scampi alla griglia. 
				
				Salmone \\
				Capparossoli, \\
				Risotto \\
				Al nero di seppia. 
				
				Spaghetti \\
				All'aragosta. 
				
				Spaghetti allo scoglio \\
				Pizza mari e monti. \\
				Paste di gamberi, 
				
				Ma soprattutto \\
				Acciughe \\
				Acciughe salate. 
				
				Peoci e trote, \\
				Salmone e seppie, \\
				Anguille, granchio, \\
				Scampi, salmone. 
				
				Acciughe \\
				Acciughe salate. \\
				\vspace{\baselineskip}
				\normalsize \textit{---Superlova}
			\end{multicols}
		\end{recipeblock}
	\end{minipage}
\end{center}